\documentclass[main.tex]{subfiles}


\begin{document}

%from pagenumber to up
% \phantomsection 
% \hypertarget{toc}{}

 

 

 
   

\section{Проводник в электрическом поле}


Проводники (заряженные или незаряженные), помещенные в постоянное электрическое поле (электростатическое поле), обладают рядом общих свойств.
\begin{enumerate}
    \item \emph{Напряженность поля внутри проводника везде равна нулю.} 

    Пусть незаряженный железный куб помещен в однородное поле (рис.\,\ref{fig:p109}).

\tikzfading[name=fade right,
left color=transparent!0,
right color=transparent!100]

\tikzfading[name=fade left,
left color=transparent!100,
right color=transparent!0]

\begin{figure}[H]
\centering

    \begin{tikzpicture}[font=\small,            line cap=round,                             >={Stealth[inset=0pt,round,length=3mm, angle'=20]},vec/.style={>={Stealth[inset=0pt,round,length=3.5mm, angle'=20]}}]


%solid
\fill[lightgray, semitransparent] (0,0) rectangle++ (3,3);
\fill[NavyBlue,path fading=fade right] (0,0) rectangle++ (.5,3);
\fill[red,path fading=fade left] (3,0) rectangle++ (-.5,3);
\draw[] (0,0) rectangle++ (3,3);

    
%E
% \draw[->,Green] (-2,-.9) --++ (7,0);
% \draw[->,Green] (-2,.3) --++ (7,0);
% \draw[->,Green] (-2,1.5) --++ (7,0);
% \draw[->,Green] (-2,2.7) --++ (7,0);
% \draw[->,Green] (-2,3.9) --++ (7,0);

\draw[->,Green] (-2,1.5) --++ (7,0) node[right,black] {$E_\text{внеш}$};
\draw[->,Green,yshift=1.7cm] (-2,1.5) --++ (7,0);
\draw[->,Green,yshift=-1.7cm] (-2,1.5) --++ (7,0);

%phantom centering
\phantom{\node[left] at (-2,1.5) {$E_\text{внеш}$};}


%E_i
\draw[->,red,yshift=.85cm] (2.4,1.5) --++ (-1.8,0) node[midway,black,above] {$E_i$};
\draw[->,red,yshift=-.85cm] (2.4,1.5) --++ (-1.8,0);


%red blue



    \end{tikzpicture}
\caption{Проводник в электрическом поле}
\label{fig:p109}
\end{figure}
\nointerlineskip

    Действие внешнего поля~$E_\text{внеш}$ ведет к избытку свободных электронов в левой части куба. На левой поверхности куба \emph{наводятся} (\emph{индуцируются}) избыточные отрицательные заряды (синяя область); на правой поверхности куба тогда имеются избыточные положительные заряды (красная область). Наведенные на поверхностях проводника заряды создают собственное поле~$E_i$, направленное \emph{против} внешнего поля~$E_\text{внеш}$ (рис.\,\ref{fig:p109}). 
    
    Во всей области внутри куба поле~$E_i$ полностью компенсирует внешнее поле~$E_\text{внеш}$: результирующее (суммарное) поле~$E$ внутри проводника оказывается равным $E=E_\text{внеш}-E_i=0$. 

    \item \emph{Избыточные заряды проводника располагаются лишь на его поверхности. Заряд любой области внутри проводника равен нулю.} 

    \item \emph{Линии электрического поля входят в проводник (или выходят из него) перпендикулярно его поверхности.}

    На рис.\,\ref{fig:p110} показана примерная картина линий результирующего поля для рассмотренного примера с проводящим кубом.

\begin{figure}[H]
\centering

    \begin{tikzpicture}[font=\small,            line cap=round,                             >={Stealth[inset=0pt,round,length=3mm, angle'=20]},vec/.style={>={Stealth[inset=0pt,round,length=3.5mm, angle'=20]}}]


%solid
\fill[lightgray, semitransparent] (0,0) rectangle++ (3,3);
\fill[NavyBlue,path fading=fade right] (0,0) rectangle++ (.5,3);
\fill[red,path fading=fade left] (3,0) rectangle++ (-.5,3);
\draw[] (0,0) rectangle++ (3,3);

    
%E
% \draw[->,Green] (-2,-.9) --++ (7,0);
% \draw[->,Green] (-2,.3) --++ (7,0);
% \draw[->,Green] (-2,1.5) --++ (7,0);
% \draw[->,Green] (-2,2.7) --++ (7,0);
% \draw[->,Green] (-2,3.9) --++ (7,0);

\draw[Green] (-2,1.5) --++ (2,0);
\draw[Green] (-2,3.2) --++ (1.95,0) arc [start angle=90, end angle=0, x radius=.3cm, y radius=.2cm];
\draw[Green] (-2,-.2) --++ (1.95,0) arc [start angle=-90, end angle=0, x radius=.3cm, y radius=.2cm];

\draw[Green] (3,1.5) --++ (2,0) node[black,right] {$E$};
\draw[Green] (5,3.2) --++ (-1.95,0) arc [start angle=90, end angle=180, x radius=.3cm, y radius=.2cm];
\draw[Green] (5,-.2) --++ (-1.95,0) arc [start angle=-90, end angle=-180, x radius=.3cm, y radius=.2cm];

\path[Green,decoration={markings,mark=at position {1/2} with {\arrow{>}}},postaction={decorate}] (-.85,3.2);
\path[Green,decoration={markings,mark=at position {1/2} with {\arrow{>}}},postaction={decorate}] (-.85,1.5);
\path[Green,decoration={markings,mark=at position {1/2} with {\arrow{>}}},postaction={decorate}] (-.85,-.2);

\path[Green,decoration={markings,mark=at position {1/2} with {\arrow{>}}},postaction={decorate}] (4.15,3.2);
\path[Green,decoration={markings,mark=at position {1/2} with {\arrow{>}}},postaction={decorate}] (4.15,1.5);
\path[Green,decoration={markings,mark=at position {1/2} with {\arrow{>}}},postaction={decorate}] (4.15,-.2);


%phantom centering
\phantom{\node[left] at (-2,1.5) {$E$};}


    \end{tikzpicture}
\caption{Линии поля вокруг проводника}
\label{fig:p110}
\end{figure}
\nointerlineskip

    Из рисунка~\ref{fig:p110} видно, что любая линия поля вблизи куба направлена точно под прямым углом к его поверхности.

    \item \emph{Все точки проводника имеют одинаковый потенциал.}

    Итак, между любыми двумя точками проводника напряжение равно нулю.
\end{enumerate}

    






















 
\end{document}